\casechapter{Topology for Combinatorics}{Kneser's conjecture and the Borsuk--Ulam theorem}{Transfer}{ch:kneser}
\begin{refsection}[chapters/07-kneser.bib]

\newcommand{\kneserKG}{\mathrm{KG}}
\newcommand{\kneserSG}{\mathrm{SG}}
\newcommand{\kneserCd}{\mathrm{cd}_2}

\begin{chaptersummary}
In 1955 Martin Kneser asked whether the $k$-element subsets of an $n$-element set can be sorted into fewer than $n-2k+2$ classes so that no class contains two disjoint sets. A simple sorting achieves $n-2k+2$, and every argument that counts or averages stops at $n/k$. In 1978 Lov\'asz answered the question by moving it to a sphere: disjoint sets become antipodal points, a sorting becomes a cover, and the Borsuk--Ulam theorem forbids the cover. We follow Greene's half-page version of this transfer, reuse its dictionary to split necklaces and sandwiches fairly, and count the cost: the transfer proves existence, and its objects can be hard to find.
\end{chaptersummary}

\section{Before: twenty-three years of colorings}\label{sec:kneser:before}

In 1955 Martin Kneser published a short problem in the \emph{Jahresbericht der Deutschen Mathematiker-Vereinigung} \parencite{Kneser1955}. Sort the $k$-element subsets of an $n$-element set into classes so that no class contains two disjoint sets. It is easy to do this with $n-2k+2$ classes. Can it be done with $n-2k+1$? Kneser conjectured that it cannot.\sidenote{Kneser was a number theorist, and the question came to him from his work on quadratic forms. It appeared as \emph{Aufgabe}~360, an exercise posed to the journal's readers.}

\begin{definition}[Kneser graph]\label{def:kneser:kg}
For integers $n,k\ge1$, the \term{Kneser graph} $\kneserKG(n,k)$ has the $k$-element subsets of $[n]=\{1,\dots,n\}$ as vertices, two of them adjacent when they are disjoint. More generally, for a finite family $\cF$ of nonempty finite sets, $\kneserKG(\cF)$ has vertex set $\cF$, two members adjacent when they are disjoint.
\end{definition}
\index{Kneser graph}

The color classes of a proper coloring of $\kneserKG(n,k)$ are the \term{intersecting families}: families of $k$-sets in which every two members meet. The smallest Kneser graph that is neither complete, nor edgeless, nor a perfect matching is $\kneserKG(5,2)$, the Petersen graph.

\begin{proposition}[Kneser's coloring]\label{prop:kneser:upper}
If $k\ge1$ and $n\ge2k-1$, then $\chi(\kneserKG(n,k))\le n-2k+2$.
\end{proposition}

\begin{proof}
Give a $k$-set $A$ the color $\min A$ when $\min A\le n-2k+1$, and the color $n-2k+2$ otherwise. Two sets with the same color $i\le n-2k+1$ both contain $i$. Two sets with the color $n-2k+2$ both lie in the set $\{n-2k+2,\dots,n\}$, which has $2k-1$ elements, and two $k$-subsets of a $(2k-1)$-element set meet because $k+k>2k-1$.
\end{proof}

\begin{conjecture}[\textcite{Kneser1955}]\label{conj:kneser:kneser}
If $k\ge1$ and $n\ge2k-1$, then $\chi(\kneserKG(n,k))=n-2k+2$. Equivalently, however the $k$-subsets of $[n]$ are sorted into $n-2k+1$ classes, some class contains two disjoint sets.
\end{conjecture}
\index{Kneser's conjecture}

The upper bound took a paragraph. The lower bound took twenty-three years. Small cases yield to small tricks, and the tricks do not combine. When $k=1$ the Kneser graph is the complete graph $K_n$, so $\chi=n$. When $n=2k-1$ it has no edges, and when $n=2k$ each set is adjacent only to its complement, so the graph is a perfect matching and $\chi=2$. When $n=2k+1$, work modulo $2k+1$ and put $A_j=\{j,j+2,\dots,j+2k-2\}$. The sets $A_j$ and $A_{j+1}$ are disjoint, because their elements are $j$ plus the even and the odd numbers between $0$ and $2k-1$, respectively, and these are distinct modulo $2k+1$. So $A_0,A_1,\dots,A_{2k},A_0$ is a closed walk of odd length, the graph is not bipartite, and $\chi\ge3$. With \cref{prop:kneser:upper}, these remarks settle every case with $n\le2k+1$. The first case that needs an idea is $k=2$.

\begin{proposition}\label{prop:kneser:k2}
If $n\ge3$, then $\chi(\kneserKG(n,2))=n-2$.
\end{proposition}

\begin{proof}
The plan is to find the shapes of the color classes and then to count edges. The upper bound is \cref{prop:kneser:upper}. A color class is a set of edges of the complete graph $K_n$ in which every two edges share a vertex, and such a set lies in a star (the edges at one vertex) or in a triangle. Indeed, a class with at most one edge lies in a star; otherwise it contains two edges $ab$ and $ac$. If every edge of the class passes through $a$, then the class lies in a star. If some edge avoids $a$, then that edge meets $ab$ and $ac$ away from $a$, so it is $bc$, and every further edge meets $ab$, $ac$, and $bc$, so both of its endpoints lie in $\{a,b,c\}$.

Suppose now that a proper coloring uses $c\le n-3$ colors, that $s$ of its classes lie in stars with centers $v_1,\dots,v_s$, and that the other $t=c-s$ classes lie in triangles. The vertices other than the centers, at least $n-s\ge t+3$ of them, span at least $\binom{n-s}{2}$ edges, none of which lies in a star class, so the triangle classes cover them and $\binom{n-s}{2}\le3t$. On the other hand,
\[
\binom{n-s}{2}-3t\ \ge\ \binom{t+3}{2}-3t\ =\ \frac{t^2-t+6}{2}\ >\ 0,
\]
a contradiction.
\end{proof}

For $k=3$ the same plan meets a zoo of intersecting families (stars, the lines of a Fano plane, the $3$-sets meeting a fixed triple in at least two points), and no short list of shapes plays the role of ``star or triangle''.

\subsection{The averaging barrier}

A second strategy needs no shapes, only sizes. No intersecting family is larger than a star: by the theorem of Erd\H{o}s, Ko, and Rado \parencite{ErdosKoRado1961}, if $n\ge2k$, then an intersecting family of $k$-subsets of $[n]$ has at most $\binom{n-1}{k-1}$ members. Every color class is therefore that small, and
\[
\chi(\kneserKG(n,k))\ \ge\ \binom nk\Big/\binom{n-1}{k-1}\ =\ \frac nk .
\]
\index{Erd\H{o}s--Ko--Rado theorem}%
For the Petersen graph this gives $\chi\ge5/2$, hence $\chi\ge3$, which is the truth. For $\kneserKG(8,3)$ it gives $8/3$, hence $3$, while the conjecture asks for $4$. No averaging argument can do better, and the reason has a name.

\begin{definition}[Fractional chromatic number]\label{def:kneser:frac}
A \term{fractional coloring} of a graph $G$ assigns a weight $w_I\ge0$ to each independent set $I$ of $G$ so that the sets containing each vertex have total weight at least $1$. The \term{fractional chromatic number} $\chi_f(G)$ is the least total weight $\sum_I w_I$ of a fractional coloring.
\end{definition}
\index{chromatic number!fractional}

The classes of a proper coloring, with weight $1$ each, form a fractional coloring, so $\chi_f\le\chi$. The number $\chi_f(G)$ is the value of a linear program, and by linear programming duality (\cref{ch:lpduality}) every lower bound on it is certified by a weighting of the vertices; the counting bound above uses equal weights, and summing the covering condition over all vertices shows that it bounds $\chi_f$ as well.\sidenote{The other classical bounds are no stronger. The eigenvalues of $\kneserKG(n,k)$ are $(-1)^i\binom{n-k-i}{k-i}$ for $0\le i\le k$, so Hoffman's eigenvalue bound $1-\lambda_{\max}/\lambda_{\min}$ also equals $n/k$; and the semidefinite bound given by Lov\'asz's theta function never exceeds $\chi_f$ (compare \cref{ch:sos}).} For Kneser graphs that bound is the truth about $\chi_f$: the $n$ stars $S_i=\{A:i\in A\}$ are independent, each $k$-set lies in exactly $k$ of them, and weight $1/k$ on every star is a fractional coloring of total weight $n/k$. So $\chi_f(\kneserKG(n,k))=n/k$ for $n\ge2k$, and Kneser's conjecture lives in the gap between $\chi_f$ and $\chi$. For fixed $t=n-2k$, $\chi_f=2+t/k$ tends to $2$ while the conjecture says $\chi=t+2$; for $\kneserKG(2k+2,k)$ with $k\ge2$, every averaging argument stops at $\chi\ge3$, and the conjecture says $4$. A proof must use something that a fractional coloring does not have.

\subsection{A second problem of the same shape}

\index{necklace splitting}%
A problem from another world has the same anatomy. Two thieves steal an open necklace with $t$ kinds of beads, a positive even number of each kind, and want to divide it \term{fairly}, each thief getting half of the beads of each kind, with as few cuts of the string as possible.\sidenote{The clasp does not count as a cut. \textcite{AlonWest1986} mention an application of the problem to the design of VLSI circuits.} If the beads of each kind sit together, then every block must be cut, so $t$ cuts may be necessary (Exercise~\ref{ex:kneser:necklace}). \textcite{GoldbergWest1985} proved that $t$ cuts always suffice, by an induction on $t$ through a continuous version of the problem that used detailed topological arguments \parencite{AlonWest1986}. The continuous version had appeared earlier in approximation theory \parencite{HobbyRice1965}, with a proof that was hard mainly because of an auxiliary odd map it had to construct \parencite{Pinkus1976}. With two kinds of beads an elementary argument works.

\begin{proposition}\label{prop:kneser:twokinds}
A necklace with two kinds of beads, each occurring an even number of times, can be divided fairly with at most two cuts.
\end{proposition}

\begin{proof}
Say there are $2a$ beads of the first kind and $2b$ of the second, and put $m=2a+2b$. For $s=0,1,\dots,m/2$ let $W_s$ be the window formed by the $m/2$ consecutive beads $s+1,\dots,s+m/2$, and let $f(s)$ be the number of beads of the first kind in $W_s$. Sliding the window by one bead changes $f$ by at most $1$. The windows $W_0$ and $W_{m/2}$ are complementary, so $f(0)+f(m/2)=2a$; one of the numbers $f(0)-a$ and $f(m/2)-a$ is at most $0$ and the other at least $0$. Hence $f(s)=a$ for some $s$. The first thief takes $W_s$, which holds $a$ beads of the first kind and $(a+b)-a=b$ of the second, and the second thief takes the rest; at most two cuts separate $W_s$ from the rest.
\end{proof}

With three kinds the windows no longer suffice: the necklace in \cref{fig:kneser:necklace} has no fair division with two cuts.

The two problems share three features: an existence statement for every input with a sharp number in it, small cases that yield to tricks which do not generalize, and a symmetry that is easy to overlook. In Kneser's problem the symmetry is the relation ``disjoint''; in the window argument it is the passage from $W_0$ to its complement $W_{m/2}$, which turns the error $f(0)-a$ into its negative. The move makes this symmetry the main object.

\section{The move: disjoint sets are antipodal points}\label{sec:kneser:move}

In a six-page note published in 1978, L\'aszl\'o Lov\'asz proved Kneser's conjecture \parencite{Lovasz1978}, and his proof counted nothing. He attached to every graph $G$ its \term{neighborhood complex}, the simplicial complex whose simplices are the sets of vertices with a common neighbor; proved, with the Borsuk--Ulam theorem, that if this complex is $j$-connected (every map of a sphere of dimension at most $j$ into it extends to the ball), then $\chi(G)\ge j+3$; and showed that for $\kneserKG(n,k)$ it is $(n-2k-1)$-connected. The method works for every graph and needs homotopy theory. Imre B{\'a}r{\'a}ny's two-page proof followed in the same issue \parencite{Barany1978}; it applied the Borsuk--Ulam theorem directly to points on a sphere, placed with the help of a lemma of Gale (Exercise~\ref{ex:kneser:gale}). In 2002 Joshua Greene, then an undergraduate, removed that lemma \parencite{Greene2002}.\sidenote{Greene studied at Harvey Mudd College, where his adviser was Francis Su (see Exercise~\ref{ex:kneser:rental}). The proof, with his senior thesis, earned him the 2002 Morgan Prize for undergraduate research.} We present the move in Greene's form, because it is the humblest.
\index{neighborhood complex}

\subsection{The Borsuk--Ulam theorem}

\begin{theorem}[{Borsuk--Ulam theorem; \textcite{Borsuk1933}}]\label{thm:kneser:bu}
Let $d\ge0$ and $S^d=\{x\in\R^{d+1}:\norm{x}=1\}$. The following equivalent statements hold.
\begin{enumerate}
\item For every continuous map $f:S^d\to\R^d$ there is a point $x$ with $f(x)=f(-x)$.
\item Every continuous \emph{odd} map $g:S^d\to\R^d$, one with $g(-x)=-g(x)$ for all $x$, has a zero.
\item There is no continuous odd map from $S^d$ to $S^{d-1}$.
\end{enumerate}
\end{theorem}
\index{Borsuk--Ulam theorem}

The equivalences take a line each: (i) gives (ii), because an odd map with $g(x)=g(-x)$ satisfies $g(x)=-g(x)$; (ii) gives (i) when applied to $g(x)=f(x)-f(-x)$; (ii) gives (iii) because $S^{d-1}$ omits the origin; and (iii) gives (ii) because a zero-free odd map $g$ yields the odd map $g/\norm{g}$. For $d=1$ the theorem is the intermediate value theorem: $h(x)=f(x)-f(-x)$ takes opposite values at $x$ and $-x$, so it vanishes on each half-circle joining them. For $d=2$ it says that at every moment two antipodal points of the Earth have the same temperature and pressure; for $d\ge2$ its proofs need degree theory or homology, or a combinatorial lemma of Tucker followed by a limiting argument \parencite[Chapter~2]{Matousek2003}; at bottom they are parity arguments, since an odd map of a sphere to itself has odd degree. We take \cref{thm:kneser:bu} as given. It is where the difficulty will go.\sidenote{Ulam asked the question, and \textcite{Borsuk1933} proved the theorem. The covering form for closed sets, a case of \cref{thm:kneser:ls} below, had appeared three years earlier in a monograph on variational problems by \textcite{LyusternikShnirelmann1930}.}

\subsection{The dictionary}

Here is the move. Place the ground set as points on a sphere. A direction $x\in S^d$ determines the open hemisphere $H(x)=\{y\in S^d:\ip{x}{y}>0\}$, and opposite directions determine disjoint hemispheres. So two sets of points lying in opposite hemispheres are disjoint: disjointness, the relation that defines the Kneser graph, has become antipodality, the relation that the Borsuk--Ulam theorem is about. \Cref{tab:kneser:dictionary} spells out the rest.
\index{hemisphere}

\begin{table}[t]
\small\renewcommand{\arraystretch}{1.18}%
\begin{tabularx}{\linewidth}{@{}>{\raggedright\arraybackslash}X>{\raggedright\arraybackslash}X@{}}
\toprule
On the sphere $S^d$, with $d=n-2k+1$ & In Kneser's problem\\
\midrule
a point of a set $X\subset S^d$ in general position & an element of $[n]$\\
antipodal directions $x$ and $-x$ & disjoint $k$-sets, in $H(x)$ and in $H(-x)$\\
the open set $U_i$ of directions $x$ for which $H(x)$ holds a $k$-set of color $i$; it is antipodal-free & the color class $i$, which has no two disjoint sets\\
the closed set $C$ of directions holding no $k$-set & antipodal-free, by general position\\
a cover of $S^d$ by $d+1$ antipodal-free sets & a proper coloring with $d$ colors\\
\bottomrule
\end{tabularx}
\caption{The dictionary of the move. By \cref{thm:kneser:ls}, the cover in the last line cannot exist.}
\label{tab:kneser:dictionary}
\end{table}

Read through the dictionary, a proper coloring with $d=n-2k+1$ colors yields a cover of $S^d$ by $d+1$ antipodal-free sets (\cref{fig:kneser:hemispheres}), which the Borsuk--Ulam theorem forbids. That is the whole proof; \cref{sec:kneser:after1} writes it out.

\begin{figure}[t]
\centering
\begin{tikzpicture}[scale=1.15]
  \def\R{1.6}\def\b{0.42}
  \fill[seafoam] (\R,0) arc[start angle=0,end angle=180,radius=\R]
     arc[start angle=180,end angle=360,x radius=\R,y radius=\b] -- cycle;
  \fill[seamist] (-\R,0) arc[start angle=180,end angle=360,radius=\R]
     arc[start angle=0,end angle=-180,x radius=\R,y radius=\b] -- cycle;
  \draw[seadeep,line width=0.6pt] (0,0) circle[radius=\R];
  \draw[sea,dashed,line width=0.45pt] (\R,0) arc[start angle=0,end angle=180,x radius=\R,y radius=\b];
  \draw[sea,line width=0.9pt] (-\R,0) arc[start angle=180,end angle=360,x radius=\R,y radius=\b];
  \fill[seadeep] (0,1.544) circle[radius=0.045];
  \node[inner sep=1pt] at (0.2,1.33) {$x$};
  \draw[muted,line width=0.5pt] (0,-1.544) circle[radius=0.045];
  \node[inner sep=1pt,text=muted] at (-0.3,-1.36) {$-x$};
  \foreach \p in {(-0.776,0.690),(0.500,1.107),(-0.497,-1.186),(0.659,-1.409),(-1.357,-0.223),(-0.056,-0.420),(1.357,-0.223)}
     \fill[seadeep] \p circle[radius=0.06];
  \node at (0.05,0.80) {$H(x)$};
  \node at (0.42,-0.92) {$H(-x)$};
  \node[right] at (\R+0.06,-0.12) {$S^2\cap x^{\perp}$};
\end{tikzpicture}
\caption{Greene's argument on $S^2$, for $n=7$ and $k=3$, so that $d=2$. If neither open hemisphere $H(x)$ nor $H(-x)$ contains three points of $X$, then at least $3=d+1$ points lie on the great circle $S^2\cap x^\perp$, which general position forbids. Only the near side is drawn.}
\label{fig:kneser:hemispheres}
\end{figure}

Why is the move natural in hindsight? The averaging barrier says that a proof must use something a fractional coloring cannot satisfy, and a parity statement such as the Borsuk--Ulam theorem is of that kind. The symmetry was in plain sight: disjointness is a symmetric relation, and its instances come in opposite pairs, like the windows $W_0$ and $W_{m/2}$. And Kneser's number counts dimensions: if both open hemispheres of a direction $x$ hold fewer than $k$ of the $n$ points, then at least $n-2(k-1)=d+1$ points lie on the great sphere $S^d\cap x^\perp$, one more than general position allows. What one had to notice is that the $k$-sets need no elaborate representation. This is Grothendieck's advice to stoop low: the space is the humblest one with a free symmetry, a sphere with its antipodal map, and the sets $U_i$ are the most naive sets that a coloring defines on it.

The pattern recurs throughout topological combinatorics, where it is often called the configuration space/test map scheme: build from the problem a space with a fixed-point-free involution, turn a hypothetical counterexample into an equivariant map or an antipodal-free cover of that space, and show by topology that none exists. \Cref{sec:kneser:after2} runs the same three steps for necklaces and sandwiches, where the involution swaps the two thieves or the two sides of a hyperplane.
\index{configuration space}

\section{After: Kneser's conjecture in half a page}\label{sec:kneser:after1}

\subsection{The Borsuk--Ulam theorem as a covering theorem}

Call a set $A\subseteq S^d$ \term{antipodal-free} if it contains no pair of points $x$ and $-x$. Lyusternik and Shnirelmann proved the following theorem for closed sets; Greene's proof needs the mixed form.
\index{antipodal-free set}

\begin{theorem}[Lyusternik--Shnirelmann]\label{thm:kneser:ls}
Let $d\ge0$. If $S^d$ is covered by $d+1$ sets, each of which is open or closed, then one of the sets contains a pair of antipodal points.
\end{theorem}
\index{Lyusternik--Shnirelmann theorem}

\begin{proof}
The plan is to prove the case of closed sets from \cref{thm:kneser:bu}, and then to reduce the general case to it by thickening the closed sets and shrinking the open ones.

\emph{Closed sets.} Let $A_1,\dots,A_{d+1}$ be closed and cover $S^d$. Define $f:S^d\to\R^d$ by $f_i(x)=\dist(x,A_i)$ for $i\le d$, with the convention $\dist(x,\emptyset)=1$. The map $f$ is continuous, so \cref{thm:kneser:bu} gives a point $x$ with $f(x)=f(-x)$. If $f_i(x)=0$ for some $i\le d$, then $x$ and $-x$ both lie in the closed set $A_i$. Otherwise neither $x$ nor $-x$ lies in $A_1\cup\dots\cup A_d$, so both lie in $A_{d+1}$.

\emph{The general case.} Suppose, for a contradiction, that $S^d$ is covered by antipodal-free sets $A_1,\dots,A_{d+1}$, each open or closed. If $A_i$ is closed and nonempty, then the compact sets $A_i$ and $-A_i$ are disjoint, so $\delta_i=\dist(A_i,-A_i)>0$; replace $A_i$ by its open neighborhood $\{y:\dist(y,A_i)<\delta_i/2\}$. The new set is still antipodal-free: if $y$ and $-y$ were within $\delta_i/2$ of points $a$ and $b$ of $A_i$, then $\norm{a-(-b)}\le\norm{a-y}+\norm{y+b}<\delta_i$, although $-b\in-A_i$. Now all the sets are open (the empty set is open), antipodal-free, and therefore different from $S^d$. The continuous function $g(y)=\max_i\dist(y,S^d\setminus A_i)$ is positive everywhere, because each $y$ lies in some open $A_i$; by compactness $g\ge\eps$ for some $\eps>0$. The closed sets $F_i=\{y:\dist(y,S^d\setminus A_i)\ge\eps\}$ cover $S^d$ and satisfy $F_i\subseteq A_i$, so they are antipodal-free, which contradicts the case of closed sets.
\end{proof}

\subsection{Greene's proof}

Greene's proof needs one lemma about points and one about colorings.

\begin{lemma}[General position]\label{lem:kneser:gp}
For all $n,d\ge1$ there is a set of $n$ points on $S^d$ no $d+1$ of which lie on a common great sphere $S^d\cap h$, $h$ a linear hyperplane of $\R^{d+1}$.
\end{lemma}
\index{general position}

\begin{proof}
Let $\gamma(s)=(1,s,s^2,\dots,s^d)\in\R^{d+1}$ and take the points $\gamma(i)/\norm{\gamma(i)}$ for $i=1,\dots,n$. Points on a common great sphere lie in a hyperplane through the origin, so they are linearly dependent. But every $d+1$ of the vectors $\gamma(i)$ are linearly independent, because the Vandermonde determinant $\prod_{p<q}(i_q-i_p)$ of $\gamma(i_1),\dots,\gamma(i_{d+1})$ is not zero, and rescaling by positive numbers preserves independence.
\end{proof}

\begin{lemma}[Hemispheres]\label{lem:kneser:hemi}
Let $d\ge1$ and let $\cF$ be a finite family of nonempty finite subsets of $S^d$. If $\kneserKG(\cF)$ has a proper coloring with $d$ colors, then there is a point $x\in S^d$ such that neither $H(x)$ nor $H(-x)$ contains a member of $\cF$.
\end{lemma}

\begin{proof}
Fix such a coloring. For each color $i$ let $U_i$ be the set of directions $x$ for which $H(x)$ contains a member of $\cF$ of color $i$, and let $C=S^d\setminus(U_1\cup\dots\cup U_d)$. Each $U_i$ is open, because for a finite set $A$ the condition $A\subseteq H(x)$ is a finite list of strict inequalities $\ip{x}{a}>0$; so $C$ is closed. The $d+1$ sets $U_1,\dots,U_d,C$ cover $S^d$, so by \cref{thm:kneser:ls} one of them contains an antipodal pair $x,-x$. That set is not a $U_i$. Otherwise $H(x)$ and $H(-x)$ would contain members $A$ and $B$ of color $i$, disjoint because the hemispheres are disjoint and distinct because they are nonempty, hence adjacent in $\kneserKG(\cF)$ and of the same color. So $x$ and $-x$ lie in $C$, which says that neither $H(x)$ nor $H(-x)$ contains a member of $\cF$.
\end{proof}

\begin{theorem}[{Kneser's conjecture; \textcite{Lovasz1978}}]\label{thm:kneser:kneser}
If $n\ge2k\ge2$, then $\chi(\kneserKG(n,k))=n-2k+2$.
\end{theorem}

\begin{proof}[Proof \parencite{Greene2002}]
The upper bound is \cref{prop:kneser:upper}. For the lower bound, put $d=n-2k+1\ge1$ and suppose, for a contradiction, that $\kneserKG(n,k)$ has a proper coloring with $d$ colors. Identify $[n]$ with a set $X$ of $n$ points on $S^d$ in general position (\cref{lem:kneser:gp}), and apply \cref{lem:kneser:hemi} to the family of $k$-subsets of $X$. It gives a direction $x$ such that each of $H(x)$ and $H(-x)$ contains at most $k-1$ points of $X$. The remaining $n-2(k-1)=d+1$ or more points of $X$ lie on the great sphere $S^d\cap x^\perp$, which general position forbids.
\end{proof}

With its two lemmas, this proof is no longer than the proof of \cref{prop:kneser:k2}, which handles $k=2$ alone, and it is short for a reason: it never asks what an intersecting family looks like. The coloring enters only through the sets $U_i$, and the size $k$ only in the last line. Lov\'asz needed a connectivity theorem and B{\'a}r{\'a}ny a lemma of Gale; Greene needs a Vandermonde determinant.

\subsection{The same proof, read again}

The size $k$ was used only in the last line. For an arbitrary set system, that line asks how many points must go before the rest splits into two parts, neither containing a member of the system.

\begin{definition}[Colorability defect]\label{def:kneser:cd}
Let $\cF$ be a family of nonempty subsets of a finite set $X$. Its \term{$2$-colorability defect} $\kneserCd(\cF)$ is the least size of a set $Y\subseteq X$ such that $X\setminus Y$ can be split into two parts, neither of which contains a member of $\cF$.
\end{definition}
\index{colorability defect}

\begin{theorem}[\textcite{Dolnikov1988}]\label{thm:kneser:dolnikov}
Let $\cF$ be a family of nonempty subsets of a finite set $X$. Then $\chi(\kneserKG(\cF))\ge\kneserCd(\cF)$.
\end{theorem}
\index{Dol'nikov's theorem}

\begin{proof}
Put $d=\chi(\kneserKG(\cF))$. If $d=0$, then $\cF$ is empty and $\kneserCd(\cF)=0$. Otherwise identify $X$ with a set of points on $S^d$ in general position (\cref{lem:kneser:gp}), and apply \cref{lem:kneser:hemi} to a proper coloring with $d$ colors. It gives a direction $x$ such that no member of $\cF$ lies in $H(x)$ or in $H(-x)$. Let $Y$ be the set of points of $X$ on the great sphere $S^d\cap x^\perp$; general position gives $|Y|\le d$. The parts $X\cap H(x)$ and $X\cap H(-x)$ split $X\setminus Y$, and neither contains a member of $\cF$. Hence $\kneserCd(\cF)\le|Y|\le d$.
\end{proof}

\Cref{thm:kneser:dolnikov} contains \cref{thm:kneser:kneser}. Let $n\ge2k$ and let $\cF$ consist of the $k$-subsets of $[n]$. Removing $n-2k+2$ elements leaves $2k-2$, which split into two halves of size $k-1$; removing fewer leaves at least $2k-1$ elements, and one of two parts then has $k$ of them. So $\kneserCd(\cF)=n-2k+2$. For $k=2$ the computation takes one line and replaces the whole proof of \cref{prop:kneser:k2}: $\cF$ is the edge set of $K_n$, and a complete graph splits into two parts without edges only when it has at most two vertices, so $\kneserCd(\cF)=n-2$.

Dol'nikov's theorem bounds the chromatic number of every finite graph $G$, because $G$ is a Kneser graph $\kneserKG(\cF)$: let $F_v$ consist of $v$ itself and the non-edges of $G$ at $v$, so that $F_u\cap F_v$ is empty exactly when $uv$ is an edge.\sidenote{\textcite{Ziegler2002} later gave a combinatorial proof of Dol'nikov's theorem and of its hypergraph versions.}

Greene's dictionary forgets the order of $[n]$, and something is lost. Call a $k$-subset of $[n]$ \term{stable} if it contains no two elements that are consecutive modulo $n$.

\begin{theorem}[\textcite{Schrijver1978}]\label{thm:kneser:schrijver}
Let $n\ge2k\ge2$. The \term{stable Kneser graph} $\kneserSG(n,k)$, induced in $\kneserKG(n,k)$ by the stable $k$-sets, has chromatic number $n-2k+2$, and deleting a vertex lowers its chromatic number.
\end{theorem}
\index{stable Kneser graph}

Greene's argument does not prove this: its hemispheres contain $k$-sets but not necessarily stable ones, and Dol'nikov's bound for stable sets is only $\max\{0,n-4k+4\}$ (Exercise~\ref{ex:kneser:dolnikov}). A proof in this spirit places the points with more care; \textcite{Ziegler2002} gave a combinatorial one. The move also reaches hypergraphs. The \term{Kneser hypergraph} $\kneserKG^r(n,k)$ has the $k$-subsets of $[n]$ as vertices and the $r$-tuples of pairwise disjoint $k$-sets as edges, and a proper coloring leaves no edge monochromatic.

\begin{theorem}[\textcite{AlonFranklLovasz1986}]\label{thm:kneser:afl}
If $r\ge2$ and $n\ge rk$, then $\chi(\kneserKG^r(n,k))=\ceil{(n-r(k-1))/(r-1)}$.
\end{theorem}
\index{Kneser hypergraph}

For $r=2$ this is \cref{thm:kneser:kneser}. Erd\H{o}s had conjectured the formula in 1973. The proof is again topological, with symmetries of prime order in place of the antipodal map.

\subsection{Combinatorics follows}

Once topology had found the reason, combinatorics could follow it. Tucker's lemma is a discrete form of the Borsuk--Ulam theorem \parencite{Tucker1946}; we need the special case that \textcite{Matousek2004} used. For $x,y\in\{+,-,0\}^n$ write $x\preceq y$ when $y$ agrees with $x$ wherever $x$ is nonzero, and write $-x$ for $x$ with all signs reversed.

\begin{theorem}[Tucker's lemma, octahedral form]\label{thm:kneser:tucker}
Let $n,m\ge1$, and let $\lambda$ assign to each nonzero $x\in\{+,-,0\}^n$ a label $\lambda(x)\in\{\pm1,\dots,\pm m\}$ with $\lambda(-x)=-\lambda(x)$. If no two sign vectors $x\preceq y$ satisfy $\lambda(x)=-\lambda(y)$, then $m\ge n$.
\end{theorem}
\index{Tucker's lemma}

The nonzero sign vectors are the vertices of a triangulated copy of $S^{n-1}$, and a labeling without complementary pairs defines an odd map into $S^{m-1}$, so \cref{thm:kneser:tucker} follows from \cref{thm:kneser:bu} (Exercise~\ref{ex:kneser:tucker}); it also has a direct combinatorial proof by a path-following parity argument \parencite[Chapter~2]{Matousek2003}.

\begin{theorem}[\textcite{Matousek2004}]\label{thm:kneser:matousek}
Tucker's lemma implies Kneser's conjecture: if $n\ge2k\ge2$, then $\chi(\kneserKG(n,k))\ge n-2k+2$.
\end{theorem}

\begin{proof}
The plan is to turn a proper coloring $c$ of $\kneserKG(n,k)$ with $n-2k+1$ colors into a labeling with labels $\pm1,\dots,\pm(n-1)$ and no complementary pair, which contradicts \cref{thm:kneser:tucker} with $m=n-1$. For a nonzero sign vector $x$ let $x^+=\{i:x_i=+\}$, $x^-=\{i:x_i=-\}$, and $|x|=|x^+|+|x^-|$.

\emph{Small vectors.} If $|x^+|<k$ and $|x^-|<k$, then $|x|\le2k-2$; put $\lambda(x)=\pm|x|$, with the sign of the first nonzero coordinate of $x$.

\emph{Large vectors.} Otherwise $x^+$ or $x^-$ contains a $k$-set. Let $c(x)$ be the largest color of a $k$-set contained in $x^+$ or in $x^-$. This color is attained on one side only, because a $k$-set in $x^+$ and a $k$-set in $x^-$ are disjoint, hence adjacent, hence of different colors. Put $\lambda(x)=2k-2+c(x)$ if $c(x)$ is attained in $x^+$, and $\lambda(x)=-(2k-2+c(x))$ if it is attained in $x^-$.

\emph{The check.} The labels lie in $\{\pm1,\dots,\pm(n-1)\}$, because $c(x)\le n-2k+1$. Replacing $x$ by $-x$ swaps $x^+$ and $x^-$, so $\lambda(-x)=-\lambda(x)$. Let $x\preceq y$; then $x^+\subseteq y^+$ and $x^-\subseteq y^-$, so $y$ is large whenever $x$ is. If both are small, then $|\lambda(x)|=|x|\le|y|=|\lambda(y)|$, with equality only when $x=y$. If $x$ is small and $y$ is large, then $|\lambda(x)|\le2k-2<|\lambda(y)|$. If both are large, then $c(x)\le c(y)$. When $c(x)<c(y)$ the labels differ in absolute value, and when $c(x)=c(y)$ a $k$-set of that color lies on the same side of $x$ and of $y$, so $\lambda(x)=\lambda(y)$. In no case is $\lambda(x)=-\lambda(y)$.
\end{proof}

The labels transcribe the topological proof: a sign vector plays a direction, $x^+$ and $x^-$ the two hemispheres, and the largest color the sets $U_i$. The combinatorial proof reached \emph{Combinatorica} in 2000, twenty-two years after the topological one.

\section{After: fair division on spheres}\label{sec:kneser:after2}

\subsection{Splitting necklaces}

The same move splits necklaces. The space is again a sphere, and the involution now swaps the two thieves.

\begin{theorem}[\textcite{HobbyRice1965}]\label{thm:kneser:hobbyrice}
Let $f_1,\dots,f_t$ be integrable functions on $[0,1]$. There are points $0=z_0\le z_1\le\dots\le z_{t+1}=1$ and signs $\eps_1,\dots,\eps_{t+1}\in\{-1,1\}$ such that
\[
\sum_{j=1}^{t+1}\eps_j\int_{z_{j-1}}^{z_j}f_i\ =\ 0\qquad\text{for } i=1,\dots,t.
\]
\end{theorem}
\index{Hobby--Rice theorem}

\begin{proof}[Proof \parencite{Pinkus1976,AlonWest1986}]
The plan is to let a point of $S^t$ choose both the cuts and the signs, and to apply \cref{thm:kneser:bu} to the resulting odd map. For $x\in S^t\subset\R^{t+1}$ put $z_0(x)=0$ and $z_j(x)=x_1^2+\dots+x_j^2$ for $1\le j\le t+1$, so that $z_{t+1}(x)=1$, and define $g:S^t\to\R^t$ by
\[
g_i(x)\ =\ \sum_{j=1}^{t+1}\sgn(x_j)\int_{z_{j-1}(x)}^{z_j(x)}f_i \qquad (i=1,\dots,t).
\]
Each term is continuous. Near a point with $x_j\ne0$ the sign is constant and the integral depends continuously on its endpoints. At a point with $x_j=0$ the term vanishes, and at a nearby point $y$ its absolute value is at most the integral of $\abs{f_i}$ over an interval of length $y_j^2$, which is small because the integral of $\abs{f_i}$ over every set of small measure is small. Replacing $x$ by $-x$ leaves every $z_j$ unchanged and reverses every sign, so $g$ is odd, and \cref{thm:kneser:bu} gives a point $x$ with $g(x)=0$. Take $z_j=z_j(x)$, $\eps_j=\sgn(x_j)$ when $x_j\ne0$, and $\eps_j=1$ when $x_j=0$; in the last case $z_j=z_{j-1}$ and the term vanishes, so the displayed sums are the numbers $g_i(x)=0$.
\end{proof}

\begin{figure}[t]
\centering
\begin{tikzpicture}[x=0.82cm,y=0.82cm]
  \draw[muted,line width=0.5pt] (0.45,0) -- (12.55,0);
  \foreach \i in {1,2,4,10} \filldraw[fill=sea,draw=seadeep,line width=0.4pt] (\i,0) circle[radius=0.27];
  \foreach \i in {3,5,6,8} \filldraw[fill=sandline,draw=seadeep,line width=0.4pt] (\i,0) circle[radius=0.27];
  \foreach \i in {7,9,11,12} \filldraw[fill=white,draw=seadeep,line width=0.4pt] (\i,0) circle[radius=0.27];
  \foreach \c in {3.5,6.5,9.5} \draw[seadeep,line width=1.1pt] (\c,-0.55) -- (\c,0.55);
  \foreach \c/\s in {2/+,5/-,8/+,11/-} \node at (\c,0.9) {$\s$};
  \foreach \l/\r/\j in {0.6/3.4/1,3.6/6.4/2,6.6/9.4/3,9.6/12.4/4}
     \draw[muted,decorate,decoration={brace,mirror,amplitude=4pt}] (\l,-0.45) -- (\r,-0.45)
        node[midway,below=5pt,text=seadeep] {$x_{\j}^2$};
\end{tikzpicture}
\caption{A fair division with three cuts of a necklace with three kinds of beads; no division with two cuts is fair. It corresponds to the point $x=\frac12(1,-1,1,-1)$ of $S^3$ in the proof of \cref{thm:kneser:hobbyrice}: piece $j$ has length $x_j^2$ and goes to the thief named by the sign of $x_j$, and $-x$ swaps the thieves.}
\label{fig:kneser:necklace}
\end{figure}

To pass from lengths to beads, place the $m$ beads of an open necklace on $[0,m]$, the $j$th bead occupying $[j-1,j]$, and let $a_i$ be the number of beads of kind $i$, which we assume even. A \term{division} of $[0,m]$ is a finite set of cut points in $(0,m)$ together with an assignment of each resulting piece to one of the two thieves. It is fair if, for every $i$, each thief receives a total length $a_i/2$ of the beads of kind $i$. Cuts at integers fall between beads, so a fair division with all its cuts at integers is a fair division of the necklace.\sidenote{\textcite{Pinkus1976} used the sphere of the $\ell_1$ norm and lengths $\abs{x_j}$; the squares follow \textcite{AlonWest1986}, who also gave the rounding argument below and noted that their approach was close to Pinkus's.}

\begin{lemma}[Rounding]\label{lem:kneser:rounding}
If $[0,m]$ has a fair division with at most $t$ cuts, then it has a fair division with at most $t$ cuts, all at integers.
\end{lemma}

\begin{proof}
The plan is to take an extremal fair division and to show that a cut inside a bead could be slid, contradicting extremality. Among the fair divisions with at most $t$ cuts, choose one with the fewest cuts and, among those, with the fewest cuts at non-integers; call the latter \emph{bad}. Each cut separates pieces given to different thieves, since otherwise deleting it would leave a fair division with fewer cuts. Suppose, for a contradiction, that a bad cut $c$ lies inside a bead $b$ of kind $i$. For a bead $b'$ let $\varphi(b')$ be the length of $b'$ given to the first thief. A bead with no bad cut has $\varphi(b')\in\{0,1\}$, while $0<\varphi(b)<1$ because $c$ separates the thieves. Fairness says that the numbers $\varphi(b')$ over the beads $b'$ of kind $i$ add up to the integer $a_i/2$, so a second bead $b''$ of kind $i$ contains a bad cut $c''$. Move $c$ and $c''$ at unit speed, in the directions that increase $\varphi(b)$ and decrease $\varphi(b'')$. Only the beads $b$ and $b''$ change, and the total for kind $i$ does not, so the division stays fair. Stop at the first moment at which a moving cut reaches an integer or meets another cut. If a cut reaches an integer strictly between $0$ and $m$ that carries no other cut, then the number of bad cuts has dropped and the number of cuts has not changed. Otherwise some piece has shrunk to a point. If it was an end piece, then the cut bounding it can be deleted. If it lay between two cuts, then the pieces on either side belong to the same thief (each of the two cuts separated thieves), and both cuts can be deleted. Every outcome contradicts the choice of the division.
\end{proof}

\begin{theorem}[{Necklace splitting; \textcite{GoldbergWest1985}, \textcite{AlonWest1986}}]\label{thm:kneser:necklace}
An open necklace with $t$ kinds of beads, each kind occurring an even number of times, can be divided fairly between two thieves with at most $t$ cuts.
\end{theorem}

\begin{proof}
Rescale $[0,m]$ to $[0,1]$ and apply \cref{thm:kneser:hobbyrice} to the indicator functions $f_i$ of the beads of kind $i$. Cutting at the distinct points among $z_1,\dots,z_t$ that lie strictly inside the interval, and giving the pieces with $\eps_j=1$ to the first thief, yields a fair division with at most $t$ cuts, because each thief receives half the length of each kind. \Cref{lem:kneser:rounding} moves the cuts between beads.
\end{proof}

Goldberg and West needed an induction on $t$ with detailed topology. Here the topology is one application of \cref{thm:kneser:bu}, to a map that writes itself once the sphere is recognized as the space of divisions with at most $t$ cuts, and the antipodal map as the exchange of the thieves. The window argument of \cref{prop:kneser:twokinds} is the case $t=2$ seen from below, with complementation of windows as the antipodal map.

\begin{theorem}[\textcite{Alon1987}]\label{thm:kneser:alon}
If a necklace has $t$ kinds of beads and the number of beads of each kind is divisible by $q$, then it can be divided fairly among $q$ thieves with at most $(q-1)t$ cuts.
\end{theorem}

The bound is attained when the beads of each kind are contiguous. Splitting in stages derives the case of $q_1q_2$ thieves from the cases $q_1$ and $q_2$, with $(q_1-1)t+q_1(q_2-1)t=(q_1q_2-1)t$ cuts; for prime $q$ Alon used a generalization of the Borsuk--Ulam theorem in which a cyclic symmetry of order $q$ replaces the antipodal map.

\subsection{Ham sandwiches}

A third dictionary lets the sphere name half-spaces.

\begin{theorem}[{Ham sandwich theorem; \textcite{StoneTukey1942}}]\label{thm:kneser:ham}
Let $d\ge1$ and let $\mu_1,\dots,\mu_d$ be finite Borel measures on $\R^d$, each vanishing on every hyperplane. There is a hyperplane each of whose open half-spaces has $\mu_i$-measure $\mu_i(\R^d)/2$, for $i=1,\dots,d$.
\end{theorem}
\index{ham sandwich theorem}

\begin{proof}
The plan is to let a point $u\in S^d$ name a half-space, so that $-u$ names the opposite one, and to apply \cref{thm:kneser:bu}. Write $u=(u_0,u')$ with $u_0\in\R$ and $u'\in\R^d$, let $P(u)=\{y\in\R^d:\ip{u'}{y}>u_0\}$, and put $f_i(u)=\mu_i(P(u))$; then $P(-u)=\{y:\ip{u'}{y}<u_0\}$. Each $f_i$ is continuous. Indeed, if $u_m\to u$, then the indicator of $P(u_m)$ converges to that of $P(u)$ at every $y$ off the set $\{y:\ip{u'}{y}=u_0\}$, which is a hyperplane when $u'\ne0$ and is empty when $u'=0$, hence $\mu_i$-null; and dominated convergence applies. \Cref{thm:kneser:bu} gives a point $u$ with $f(u)=f(-u)$. If $u'=0$, then $P(u)$ and $P(-u)$ are $\R^d$ and $\emptyset$ in some order, so every $\mu_i$ is zero and every hyperplane works. If $u'\ne0$, then the hyperplane $\{y:\ip{u'}{y}=u_0\}$ is $\mu_i$-null and its two open sides have equal measure, which is therefore $\mu_i(\R^d)/2$.
\end{proof}

This is new reach: no averaging argument bisects three solids in space with one plane. Stone and Tukey proved the theorem in all dimensions in 1942; the three-dimensional case is older. A version for finite point sets follows by a limiting argument (Exercise~\ref{ex:kneser:discreteham}).

\section{Costs and limits}\label{sec:kneser:costs}

\paragraph{Where did the difficulty go?} Into \cref{thm:kneser:bu}, and nowhere else. Given the theorem, every proof in \cref{sec:kneser:after1,sec:kneser:after2} fits in a page. The theorem itself needs degree theory or homology, or Tucker's lemma and a limiting argument, and Tucker's lemma needs a parity argument. Parity is exactly what the averaging arguments of \cref{sec:kneser:before} lacked: a fractional coloring has no parity, while an odd map of a sphere does. The price is paid once. Kneser's conjecture, Dol'nikov's theorem, the necklace theorem, and the ham sandwich theorem each paid only for a dictionary (a point set in general position, a sphere of divisions, a sphere of half-spaces), and each dictionary took a paragraph. So the move is a dissolution and not a renaming: no application proves the Borsuk--Ulam theorem again. What recurs is the work of finding the dictionary, and it is real work: Pinkus's simple odd map came eleven years after Hobby and Rice's elaborate one.

\paragraph{What did not dissolve.} First, witnesses. The proofs say that a fair division exists, not where it is. Matoušek's proof can be run as a search, since Tucker's lemma has a path-following proof, but the search ranges over the $3^n-1$ nonzero sign vectors. For necklaces the obstacle is a theorem: finding a fair division between two thieves with at most $t$ cuts is PPA-complete, and so is finding a ham sandwich cut of finite point sets \parencite{FilosRatsikasGoldberg2019}.\sidenote{PPA, a class defined by Papadimitriou, consists of the search problems reducible to this one: in an exponentially large graph described by a circuit, given a vertex of odd degree, find another. One exists, since a finite graph has an even number of vertices of odd degree.} A polynomial-time algorithm for either would put all of PPA in polynomial time. Topology certifies existence here as linear programming duality certifies optimality in \cref{ch:lpduality}, but no efficient counterpart to the dynamics of \cref{ch:regret} or the resampling of \cref{ch:compression} is expected.
\index{PPA}

Second, typical graphs. Every graph is a Kneser graph $\kneserKG(\cF)$, so \cref{thm:kneser:dolnikov} bounds every chromatic number, but topological bounds are weak for most graphs. \textcite{SimonyiTardos2006} showed, with Ky Fan's generalization of the Borsuk--Ulam theorem, that when a topological bound of the kind used in this chapter proves $\chi(G)\ge t$, every proper coloring of $G$ contains a complete bipartite subgraph $K_{\ceil{t/2},\floor{t/2}}$ whose $t$ vertices receive different colors. By first-moment computations, with high probability the random graph $G(n,1/2)$ contains no $K_{s,s}$ with $s\ge2\log_2n+1$ and no independent set of $2\log_2n+2$ vertices; so such bounds give only $O(\log n)$ for it, while its chromatic number exceeds $n/(2\log_2n+2)$.

Third, lost information. Greene's dictionary forgets the cyclic order on which Schrijver's theorem depends, and Dol'nikov's bound sees only $\max\{0,n-4k+4\}$ of the $n-2k+2$ colors that stable sets need.

Fourth, more than two parts. For Kneser hypergraphs and for $q$ thieves the antipodal map gives way to cyclic symmetries of prime order, and sometimes topology then has nothing to say. Tverberg's theorem splits every $(r-1)(d+1)+1$ points of $\R^d$ into $r$ parts whose convex hulls share a point. Its topological version, for continuous maps from a simplex of dimension $(r-1)(d+1)$ to $\R^d$, was proved by equivariant methods when $r$ is a prime power, and it is false for every other $r$ in high enough dimension \parencite{Frick2015}.\sidenote{Frick combined a generalized van Kampen theorem of Mabillard and Wagner with a constraint method of Blagojevi\'c, Frick, and Ziegler; his smallest counterexample has $r=6$ and $d=19$.} When the obstruction vanishes, the transfer is silent, and the silence can mean that the statement is false.

\paragraph{Other seas.} Lov\'asz's neighborhood complex, which works for every graph, began a theory of complexes attached to graphs \parencite{Matousek2003}. \textcite{Ziegler2002} widened the combinatorial sea that Matoušek opened, with combinatorial proofs of the theorems of Dol'nikov, of Alon, Frankl, and Lov\'asz, and of Schrijver. \textcite{ChaseEtAl2024} found a local form of \cref{thm:kneser:ls}, in which, for $d\ge1$, every finite open cover of $S^d$ by antipodal-free sets has a point lying in at least $\ceil{(d+3)/2}$ of the sets, and used it in learning theory (Exercise~\ref{ex:kneser:replicable}). The algebraic sea of eigenvalues, theta functions, and linear programs is the right one for independence numbers, where it proves the Erd\H{o}s--Ko--Rado theorem; for Kneser's conjecture it sees only $\chi_f$.

\section{The move, distilled}\label{sec:kneser:distilled}

Here is the move of this chapter, in the format of the book.

\begin{move}{Transfer}
\movepart{In one sentence} When a combinatorial statement resists every averaging argument, find its hidden symmetry, rebuild the problem on a space where that symmetry acts freely, and let a topological obstruction forbid the counterexample.
\movepart{Triggers} A fractional relaxation far from the integral truth; objects in natural opposite pairs (disjoint sets, complementary windows, two thieves, two sides of a hyperplane); an existence claim for every input, with a sharp number that looks like a dimension; small cases that need unrelated tricks.
\movepart{Recipe} (1) Name the involution. (2) Choose the simplest space on which it acts freely, usually a sphere with its antipodal map. (3) Turn a counterexample into an equivariant map, or an antipodal-free cover, of that space. (4) Rule it out with the Borsuk--Ulam theorem in a suitable form. (5) Read the conclusion back; for a combinatorial proof, use Tucker's lemma instead.
\movepart{Anti-pattern} A target space whose obstruction is as hard to establish as the original problem; topology where averaging suffices; or a dictionary under which the equivariant map exists, so that the obstruction vanishes and the transfer proves nothing.
\end{move}

\dissolutiontest{Kneser's conjecture for all $n$ and $k$ in half a page from \cref{thm:kneser:bu}, against a classification for $k=2$ alone; necklaces in half a page, against an induction on $t$ with detailed topology.}{The size $k$ disappears (\cref{thm:kneser:dolnikov}); beads become integrable densities (\cref{thm:kneser:hobbyrice}).}{Coloring bounds, necklaces, and ham sandwiches become one statement about odd maps of spheres.}{Dol'nikov's and Schrijver's theorems, Kneser hypergraphs, many thieves, ham sandwiches in every dimension; PPA-completeness and replicability lower bounds.}

The move is a transfer: the problem is moved rather than enlarged, and the reward is an obstruction that exists only in the new setting. \Cref{ch:dimension} makes the same kind of move toward linear algebra, where the obstruction is a dimension rather than a parity. \Cref{ch:lpduality} supplies the counterpoint, since linear programming duality certifies exactly the fractional chromatic number, and \cref{ch:fourier} shows what diagonalizing a symmetry achieves when the conclusion is an identity rather than an existence statement; the free symmetry used here is a cousin of the exchangeability of \cref{ch:conformal}. The price of topological existence, PPA-hardness, contrasts with the constructive moves of \cref{ch:regret,ch:compression}, and Exercise~\ref{ex:kneser:replicable} connects this chapter to the stability notions of \cref{ch:littlestone}.

\begin{exercises}
\item \consolidation\label{ex:kneser:petersen} (a) Show that $\kneserKG(5,2)$ is $3$-regular with $10$ vertices and has no cycles of length $3$ or $4$. (b) Compute its chromatic number, its independence number, and its fractional chromatic number directly. (c) Show that the stable Kneser graph $\kneserSG(2k+1,k)$ is a cycle of length $2k+1$, so that $\chi(\kneserSG(2k+1,k))=3$ although it has only $2k+1$ of the $\binom{2k+1}{k}$ vertices of $\kneserKG(2k+1,k)$.

\item \consolidation\label{ex:kneser:dolnikov} (a) For a finite graph $H$, show that $\kneserCd(E(H))$ is the least number of vertices whose deletion leaves a bipartite graph; so \cref{thm:kneser:dolnikov} bounds the chromatic number of the complement of the line graph of $H$. (b) For $H=K_n$, recover \cref{prop:kneser:k2}. (c) For $H=C_{2m+1}$ with $m\ge2$, show that the bound is $1$ while the chromatic number is $m+1$. (d) Let $k\ge1$ and $n\ge4k-4$. Show that the family of stable $k$-subsets of $[n]$ has $\kneserCd=n-4k+4$, and conclude that \cref{thm:kneser:dolnikov}, applied to this family, does not prove \cref{thm:kneser:schrijver} when $k\ge2$.

\item \consolidation\label{ex:kneser:lsbu} (a) Show that $S^d$ can be covered by $d+2$ closed antipodal-free sets, by projecting radially the facets of a $(d+1)$-simplex that contains the origin in its interior; so \cref{thm:kneser:ls} is sharp. (b) For $d\ge1$, deduce statement (iii) of \cref{thm:kneser:bu} from the case of \cref{thm:kneser:ls} in which all sets are closed.

\item \consolidation\hard\label{ex:kneser:tucker} Prove \cref{thm:kneser:tucker} from \cref{thm:kneser:bu}. \hint{Send the sign vector $x$, with $\pm$ read as $\pm1$, to the barycenter $\sum_i x_ie_i/\abs{\{i:x_i\ne0\}}$ of a face of the cross-polytope $\conv\{\pm e_1,\dots,\pm e_n\}$; the chains $x^1\prec\dots\prec x^j$ span the simplices of a triangulation of its boundary. Send the vertex $x$ to $\sgn(\lambda(x))\,e_{\abs{\lambda(x)}}\in\R^m$, extend linearly on each simplex, and check that the image avoids $0$.}

\item \consolidation\hard\label{ex:kneser:gale} (B{\'a}r{\'a}ny's proof.) (a) Let $d\ge0$, $k\ge1$, $n=2k+d$, and $\gamma(s)=(1,s,\dots,s^d)$. Show that every open hemisphere of $S^d$ contains the point $v_i=(-1)^i\gamma(i)/\norm{\gamma(i)}$ for at least $k$ of the indices $i=1,\dots,n$. This is Gale's lemma. (b) Deduce the lower bound of \cref{thm:kneser:kneser} from the case of \cref{thm:kneser:ls} in which all sets are open, working on $S^{n-2k}$ with $n-2k+1$ colors. \hint{For (a), tilt the hemisphere $H(a)$ so that no $v_i$ lies on its boundary, without changing which points it contains. Then the numbers $\ip{a}{\gamma(i)}$, $i=1,\dots,n$, change sign at most $d$ times, being values of a nonzero polynomial of degree at most $d$, and within each of the at most $d+1$ runs of constant sign the points $v_i$ alternate between the two sides.}

\item \consolidation\label{ex:kneser:necklace} (a) Show that if each of the $t$ kinds occurs and the beads of each kind are contiguous, then every fair division uses at least $t$ cuts. (b) A closed necklace is a circle of beads. Show that a closed necklace with $t$ kinds of beads, each occurring a positive even number of times, can be divided fairly between two thieves with at most $2\ceil{t/2}$ cuts, and that this is sharp. (c) Check that the necklace of \cref{fig:kneser:necklace} has no fair division with two cuts.

\item \transfer\hard\label{ex:kneser:discreteham} Let $P_1,\dots,P_d$ be finite sets in $\R^d$. Show that some hyperplane has, for each $i$, at most $\floor{|P_i|/2}$ points of $P_i$ in each of its two open half-spaces. \hint{Replace each point by the uniform probability measure on a small ball around it, apply \cref{thm:kneser:ham}, and let the radius tend to $0$; the bisecting hyperplanes all meet a fixed ball, so a subsequence converges to a hyperplane.}

\item \transfer\hard\label{ex:kneser:radon} Let $\Delta=\conv\{e_1,\dots,e_{d+2}\}\subset\R^{d+2}$ and let $f:\Delta\to\R^d$ be continuous. Show that two disjoint faces of $\Delta$ have intersecting images under $f$. For affine $f$ this is Radon's theorem: every $d+2$ points of $\R^d$ split into two parts whose convex hulls meet. \hint{For $x\in S^{d+1}$ write $x=x^+-x^-$ with $x^+,x^-\ge0$ coordinatewise and with disjoint supports, put $s=\sum_ix^+_i$ and $r=\sum_ix^-_i$, and apply \cref{thm:kneser:bu} to the odd map $x\mapsto\bigl(s-r,\ s\,f(x^+/s)-r\,f(x^-/r)\bigr)\in\R^{d+1}$, in which $s\,f(x^+/s)$ means $0$ when $s=0$, and similarly for $r$.}

\item \transfer\harder\label{ex:kneser:rental} Three roommates rent a three-room house for a total rent of $1$. A division of the rent is a point $p$ of $\Delta=\{p\in\R^3:p_i\ge0,\ p_1+p_2+p_3=1\}$, with $p_i$ the rent of room $i$; at each $p$ each roommate accepts a nonempty set of rooms. Assume that acceptance is closed: a roommate who accepts room $i$ at each point of a sequence converging to $p$ accepts room $i$ at $p$. Assume also that the roommates are hungry: when some room is free at $p$, each roommate accepts only free rooms at $p$. Show that for some division of the rent the three roommates accept three different rooms \parencite{Su1999}. \hint{Triangulate $\Delta$ finely and pass to the barycentric subdivision; give the vertex at the barycenter of a face of dimension $j$ to roommate $j+1$, and label it by a room its owner accepts there. On the open side where $p_i=0$ the label is $i$. Check that exactly one boundary edge has labels $\{1,2\}$, count small triangles through their $\{1,2\}$-edges as in Sperner's lemma, and finish by compactness.}

\item \research\harder\label{ex:kneser:replicable} An estimator sees $N$ independent samples of a vector of $d$ independent coin flips with unknown biases $p\in[0,1]^d$ and, using internal randomness as well, outputs a vector in $\R^d$. Call it $(\eps,L,\delta)$-list-replicable if for every $p$ some list of at most $L$ vectors, each within $\eps$ of $p$ in the maximum norm, contains the output with probability at least $1-\delta$. (a) Show that if $\eps<1/2$ and $\delta<1/(L+1)$, then $L\ge d+1$. (b) \textcite{ChaseEtAl2024} used their local form of \cref{thm:kneser:ls} to bound list replicability in PAC learning; for instance, the $m$ coordinate functions on $\{0,1\}^m$ have list-replicability number $1+\floor{m/2}$, the most a class of size $m$ can have. They ask whether, for finite classes, this number is bounded by a function of the VC and dual VC dimensions. Investigate this open question, or find another stability lower bound in learning that is a Borsuk--Ulam or Sperner obstruction (compare \cref{ch:littlestone}). \hint{For (a), let $U_v$ be the set of $p$ at which the output equals $v$ with probability greater than $\delta$. Each $U_v$ is relatively open, since the probability is a polynomial in $p$; it lies within $\eps$ of $v$; and finitely many of the sets $U_v$ cover the cube. Use Lebesgue's covering theorem: a finite cover of $[0,1]^d$ by relatively open sets none of which meets two opposite facets has a point lying in at least $d+1$ of the sets.}
\end{exercises}

\notesandrefs

Kneser's problem \parencite{Kneser1955} was solved by \textcite{Lovasz1978}, whose proof introduced the neighborhood complex and is usually taken as the start of topological combinatorics. \textcite{Barany1978} combined the Lyusternik--Shnirelmann theorem with Gale's lemma, \textcite{Schrijver1978} found the stable Kneser graphs in the same year, and \textcite{Greene2002} removed Gale's lemma. \textcite{Matousek2004} gave the first combinatorial proof, which \textcite{Ziegler2002} extended to the theorems of \textcite{Dolnikov1988}, of \textcite{AlonFranklLovasz1986}, and of Schrijver. The standard reference is \textcite{Matousek2003}: its Chapter~2 treats the forms of the Borsuk--Ulam theorem and their proofs, and its Section~3.3 treats Kneser's conjecture.

\textcite{Borsuk1933} proved the theorem three years after the covering form of \textcite{LyusternikShnirelmann1930}; \textcite{Tucker1946} found the combinatorial lemma. The continuous necklace theorem is due to \textcite{HobbyRice1965} and its simple proof to \textcite{Pinkus1976}; \textcite{GoldbergWest1985} proved the discrete theorem, \textcite{AlonWest1986} gave the Borsuk--Ulam proof and the rounding step, and \textcite{Alon1987} settled the case of $q$ thieves. The ham sandwich theorem in all dimensions is due to \textcite{StoneTukey1942}, and \textcite{Su1999} explains how Sperner's lemma divides rent and cake fairly. For the limits of the method see \textcite{FilosRatsikasGoldberg2019} on PPA-completeness, \textcite{SimonyiTardos2006} on colorful bipartite subgraphs, \textcite{Frick2015} on the topological Tverberg conjecture, and \textcite{ChaseEtAl2024} on local Borsuk--Ulam theorems and replicable learning.

\printchapterbib
\end{refsection}
